\section{Hình học hàm mất mát và bệnh lý građien}
\label{sec:gradient-pathology}

Trong phần này, hình học hàm mất mát là hình học của ánh xạ
$\boldsymbol{\theta}\mapsto J_N(\boldsymbol{\theta})$ trong không gian
tham số: độ cong, các hướng ít thay đổi, các hướng biến thiên nhanh
và quan hệ giữa các thành phần mất mát. Không đồng nhất khái niệm này
với không gian hàm chứa nghiệm của phương trình.

\subsection{Građien của tổng và sự cạnh tranh giữa các ràng buộc}

Với trọng số cố định, chia \eqref{p4:eq:loss-total} cho
$\lambda_r>0$ rồi gộp các thành phần còn lại thành $J_1,\ldots,J_M$:
\begin{equation}
 J(\boldsymbol{\theta})=J_r(\boldsymbol{\theta})+
                         \sum_{j=1}^{M}\lambda_jJ_j(\boldsymbol{\theta}).
 \label{p4:eq:compact-loss}
\end{equation}
Ở đây, $M$ là số nhóm dữ kiện hoặc ràng buộc, và các $\lambda_j$ đã
được chia cho trọng số phần dư ban đầu. Phép chuẩn hóa giữ nguyên tập
cực tiểu, nhưng đổi độ lớn građien; để giữ nguyên một bước hạ dốc còn
phải đổi chiều dài bước tương ứng.

Đặt $\mathbf{g}_r=\nabla_{\boldsymbol{\theta}}J_r$,
$\mathbf{g}_j=\nabla_{\boldsymbol{\theta}}J_j$ và
$\mathbf{g}=\mathbf{g}_r+\sum_j\lambda_j\mathbf{g}_j$.
Với bước $\eta>0$, hạ dốc cập nhật
\begin{equation}
 \boldsymbol{\theta}^{+}=\boldsymbol{\theta}-\eta\mathbf{g}.
 \label{p4:eq:gradient-step}
\end{equation}
Nếu $J_j$ khả vi, khai triển bậc nhất tại tham số hiện tại cho
\begin{equation}
 J_j(\boldsymbol{\theta}^{+})-J_j(\boldsymbol{\theta})
 =-\eta\mathbf{g}_j^T\mathbf{g}+o(\eta),
 \qquad \eta\downarrow0.
 \label{p4:eq:individual-change}
\end{equation}
Do đó thành phần $J_j$ giảm với bước đủ nhỏ nếu
$\mathbf{g}_j^T\mathbf{g}>0$, nhưng có thể tăng nếu tích vô hướng âm.
Ngoài chênh lệch độ lớn, cần xét cả góc giữa các građien. Khi hai
vectơ khác không, đặt
\begin{equation}
 c_{ij}=\frac{\mathbf{g}_i^T\mathbf{g}_j}
                    {\|\mathbf{g}_i\|_2\|\mathbf{g}_j\|_2}\in[-1,1].
 \label{p4:eq:cosine}
\end{equation}
Trong đó $\|\cdot\|_2$ là chuẩn Euclid. Giá trị $c_{ij}<0$ biểu thị
sự đối nghịch cục bộ; cân bằng độ lớn không tự động loại bỏ sự đối nghịch.

Ví dụ, tại một điểm của bài toán một tham số, nếu $g_r=10$, $g_b=-1$
và $\lambda_b=1$, thì $g=9$. Bước hạ dốc làm
$J_b(\theta^+)-J_b(\theta)=9\eta+o(\eta)>0$ dù tổng mất mát giảm ở
bậc nhất với lượng $81\eta$. Đây là ví dụ về hướng cập nhật, không
phải số liệu thực nghiệm của một mạng cụ thể.

Wang và cộng sự~\cite{wang2021} nghiên cứu các mất cân bằng građien
trong PINN và liên hệ chúng với tính cứng của động lực học huấn luyện.
Tuy nhiên, không nên suy ra rằng phần dư luôn là thành phần trội trên
mọi phương trình hoặc mọi cách khử thứ nguyên.

Một giới hạn logic cần nhấn mạnh: vì các mất mát không âm,
\begin{equation}
 J_j(\boldsymbol{\theta})\leq
                \frac{J(\boldsymbol{\theta})}{\lambda_j}.
 \label{p4:eq:component-bound}
\end{equation}
Với các trọng số cố định và dương, $J\to0$ kéo theo mọi $J_j\to0$.
Vấn đề là một giá trị tổng hữu hạn được gọi là nhỏ có thể chưa đủ
cho ràng buộc có trọng số nhỏ, và các mất mát này vẫn chỉ kiểm tra
tại các điểm đã lấy mẫu.

\subsection{Hessian của bài toán bình phương tối thiểu}

Ghép các sai lệch đã nhân hệ số $\sqrt{\lambda_j/N_j}$ thành vectơ
$\mathbf{q}(\boldsymbol{\theta})\in\mathbb{R}^{S}$, trong đó $S$ là
tổng số sai lệch vô hướng. Khi đó $J=\|\mathbf{q}\|_2^2$.
Đặt ma trận Jacobi
$\mathbf{A}=D_{\boldsymbol{\theta}}\mathbf{q}\in\mathbb{R}^{S\times P}$.
Nếu các $q_a$ khả vi hai lần theo tham số thì
\begin{align}
 \nabla J&=2\mathbf{A}^{T}\mathbf{q},\label{p4:eq:ls-gradient}\\
 \nabla^2J&=2\mathbf{A}^{T}\mathbf{A}
          +2\sum_{a=1}^{S}q_a\nabla^2q_a.
          \label{p4:eq:ls-hessian}
\end{align}
Để suy ra, lấy $\partial_{\theta_p}J=2\sum_aq_a\partial_{\theta_p}q_a$
rồi lấy tiếp đạo hàm theo $\theta_q$. Quy tắc tích tạo hai số hạng
$2\sum_a(\partial_{\theta_p}q_a)(\partial_{\theta_q}q_a)$ và
$2\sum_aq_a\partial^2_{\theta_q\theta_p}q_a$, chính là hai ma trận
ở \eqref{p4:eq:ls-hessian}.

Ma trận $2\mathbf{A}^{T}\mathbf{A}$ nửa xác định dương, nhưng số hạng
còn lại có thể không như vậy. Tại nghiệm có $\mathbf{q}=0$, số hạng
thứ hai triệt tiêu; gần nghiệm, việc bỏ nó cần được biện minh bằng
độ lớn của phần dư và các đạo hàm bậc hai. Tại một điểm dừng bất kỳ,
không được mặc nhiên coi mọi trị riêng của Hessian là không âm.

Hessian có thể suy biến do ma trận Jacobi thiếu hạng. Tuy nhiên,
đối xứng hoán vị các đơn vị ẩn là đối xứng rời rạc, tự nó không sinh
một hướng tiếp tuyến có độ cong bằng không. Hơn nữa, trị riêng bằng
không chỉ nói về xấp xỉ bậc hai: ví dụ $J(\theta)=\theta^4$ có
$J''(0)=0$ nhưng không hằng trên lân cận của $0$.

\subsection{Dòng građien và tính cứng}

Giữ trọng số cố định và xét thời gian tối ưu $s$, độc lập với thời
gian vật lý $t$. Dòng građien là hệ
\begin{equation}
 \frac{\mathrm{d}\boldsymbol{\theta}}{\mathrm{d}s}
 =-\nabla J(\boldsymbol{\theta}(s)).
 \label{p4:eq:flow}
\end{equation}
Quy tắc dây chuyền cho
\begin{equation}
 \frac{\mathrm{d}}{\mathrm{d}s}J(\boldsymbol{\theta}(s))
 =\nabla J^T\frac{\mathrm{d}\boldsymbol{\theta}}{\mathrm{d}s}
 =-\|\nabla J\|_2^2\leq0.
 \label{p4:eq:flow-energy}
\end{equation}
Vì vậy tổng mất mát giảm dọc dòng liên tục, dù từng thành phần có thể
tăng như \eqref{p4:eq:individual-change} đã chỉ ra.

Quanh một cực tiểu $\boldsymbol{\theta}^{*}$ của hàm $C^2$, đặt
$\mathbf{h}=\boldsymbol{\theta}-\boldsymbol{\theta}^{*}$ và
$\mathbf{H}=\nabla^2J(\boldsymbol{\theta}^{*})\succeq0$.
Ta có $\nabla J(\boldsymbol{\theta}^{*}+\mathbf{h})
=\mathbf{H}\mathbf{h}+o(\|\mathbf{h}\|)$, nên hệ tuyến tính hóa là
$\dot{\mathbf{h}}=-\mathbf{H}\mathbf{h}$.
Nếu $\mathbf{H}\mathbf{v}_j=\mu_j\mathbf{v}_j$ với hệ vectơ riêng
trực chuẩn, thì hệ số $a_j=\mathbf{v}_j^T\mathbf{h}$ thỏa
\begin{equation}
 \dot a_j=-\mu_ja_j,
 \qquad a_j(s)=a_j(0)\exp(-\mu_js).
 \label{p4:eq:flow-modes}
\end{equation}
Trong đó, $\mu_j\geq0$ là trị riêng và $a_j$ là sai lệch theo hướng
riêng đó. Các trị riêng dương lớn và nhỏ tạo những thang thời gian
rất khác nhau; đó là cơ chế tính cứng trong mô hình tuyến tính hóa.

Hạ dốc là lược đồ Euler hiện cho \eqref{p4:eq:flow}. Trên mỗi hướng
có trị riêng dương, hệ tuyến tính cho
\begin{equation}
 a_j^{k+1}=(1-\eta\mu_j)a_j^k,
 \qquad |1-\eta\mu_j|<1
 \ \Longleftrightarrow\ 0<\eta<\frac{2}{\mu_j}.
 \label{p4:eq:euler-stability}
\end{equation}
Nếu $\mu_{\max}>0$ và $\mu_{\min}^{+}>0$ lần lượt là trị riêng lớn
nhất và trị riêng dương nhỏ nhất, đặt
$\kappa_{\mathrm{eff}}=\mu_{\max}/\mu_{\min}^{+}$.
Với $\eta=1/\mu_{\max}$, hướng chậm nhất có hệ số co
$1-1/\kappa_{\mathrm{eff}}$; sau $k$ bước, nó bị chặn bởi
$\exp(-k/\kappa_{\mathrm{eff}})$ lần sai lệch đầu. Vì vậy việc giảm
sai lệch đi hệ số $\epsilon\in(0,1)$ được bảo đảm trong mô hình này
nếu $k\geq\kappa_{\mathrm{eff}}\log(1/\epsilon)$.
Các hướng có trị riêng bằng không không co trong hệ tuyến tính;
động lực học thực ở những hướng đó có thể phụ thuộc số hạng bậc cao.
Đây là phân tích cục bộ với Hessian đóng băng, không phải tốc độ hội
tụ toàn cục của một PINN phi tuyến.

\subsection{Điều chỉnh trọng số bằng thống kê građien}

Wang và cộng sự đề xuất điều chỉnh các hệ số $\lambda_j$ dựa trên
građien từng thành phần~\cite{wang2021}. Công thức dưới đây được đối
chiếu trực tiếp với Thuật toán~1 của bản tác giả~\cite{wang2020}.
Tại bước $k$, đặt
\begin{equation}
 A_k=\max_{1\leq p\leq P}|(\mathbf{g}_{r,k})_p|,
 \qquad B_{j,k}=\frac{1}{P}\sum_{p=1}^{P}|(\mathbf{g}_{j,k})_p|.
 \label{p4:eq:gradient-statistics}
\end{equation}
Trong đó, $\mathbf{g}_{r,k}=\nabla J_r(\boldsymbol{\theta}_k)$,
$\mathbf{g}_{j,k}=\nabla J_j(\boldsymbol{\theta}_k)$; chỉ số $p$ chạy
qua các tham số, không chạy qua các điểm không gian. Tử số là độ lớn
thành phần lớn nhất, còn mẫu số là trung bình độ lớn các thành phần.
Khi $B_{j,k}>0$, chọn
\begin{equation}
 \widehat\lambda_{j,k}=\frac{A_k}{B_{j,k}},
 \qquad
 \lambda_{j,k+1}=\rho\lambda_{j,k}+(1-\rho)\widehat\lambda_{j,k},
 \quad 0<\rho<1.
 \label{p4:eq:weight-update}
\end{equation}
Ở đây, $\widehat\lambda_{j,k}$ là trọng số đề xuất và $\rho$ là hệ số
giữ lại giá trị cũ. Khi bỏ làm trơn, công thức cho đúng đẳng thức
\begin{equation}
 \frac{1}{P}\sum_{p=1}^{P}
       |\widehat\lambda_{j,k}(\mathbf{g}_{j,k})_p|=A_k.
 \label{p4:eq:balance-meaning}
\end{equation}
Do đó quy tắc cân bằng hai thống kê đã chọn, không cân bằng từng tọa
độ và không bảo đảm các građien cùng hướng.

Với quy ước của \eqref{p4:eq:weight-update}, $\rho$ gần một làm thay
đổi chậm. Nếu viết $(1-\alpha)\lambda+\alpha\widehat\lambda$ thì
$\alpha=1-\rho$. Phải giữ đúng quy ước khi so sánh các giá trị tham
số giữa tài liệu và cài đặt. Quy tắc này điều chỉnh trọng số của các
thành phần; chiều dài bước chung $\eta_k$ vẫn là một đại lượng khác.

\paragraph{Xử lý mẫu số nhỏ.}
Một biến thể có bảo vệ số học là
\begin{equation}
 \widehat\lambda_{j,k}^{\mathrm{reg}}
 =\min\left\{\lambda_{\max},\,
       \max\left\{\lambda_{\min},\,
          \frac{A_k}{B_{j,k}+\varepsilon_g}\right\}\right\},
 \label{p4:eq:regularized-weight}
\end{equation}
với $\varepsilon_g>0$ và $0<\lambda_{\min}\leq\lambda_{\max}<\infty$.
Đây là phần bổ sung để cài đặt, không phải công thức nguyên dạng của
thuật toán được dẫn. Hằng số $\varepsilon_g$ tránh chia cho không;
hai cận hạn chế sự tăng hoặc giảm quá mức. Khi dùng biến thể này,
đẳng thức \eqref{p4:eq:balance-meaning} chỉ còn là mục tiêu gần đúng.
Nếu $\mathbf{g}_{j,k}=0$ nhưng $J_j>0$, tăng trọng số không tạo được
hướng cải thiện từ một građien bằng không; cần kiểm tra tham số hóa
hoặc điểm dừng của thành phần đó.

\begin{samepage}
\noindent\textbf{Quy trình huấn luyện có điều chỉnh trọng số.}
\begin{enumerate}[leftmargin=*,label=\arabic*.]
 \item Chọn $\boldsymbol{\theta}_0$, $\lambda_{j,0}>0$, chu kỳ cập nhật
 $K\geq1$, hệ số $\rho$ và dãy bước $\eta_k>0$.
 \item Tại bước $k\equiv0\pmod K$, tính riêng các građien tại
 cùng $\boldsymbol{\theta}_k$, rồi dùng
 \eqref{p4:eq:gradient-statistics}--\eqref{p4:eq:weight-update}; nếu cần,
 thay tỉ số bằng \eqref{p4:eq:regularized-weight}. Ở các bước khác,
 giữ $\lambda_{j,k+1}=\lambda_{j,k}$.
 \item Giữ các trọng số mới cố định trong phép lấy đạo hàm và cập nhật
 \[
  \boldsymbol{\theta}_{k+1}=\boldsymbol{\theta}_k-
  \eta_k\left[\mathbf{g}_{r,k}+\sum_{j=1}^{M}
                         \lambda_{j,k+1}\mathbf{g}_{j,k}\right].
 \]
 \item Theo dõi từng $J_j$, các thống kê građien và phần dư trên tập
 kiểm tra độc lập; không chỉ theo dõi tổng mất mát có trọng số thay đổi.
\end{enumerate}
\end{samepage}

Trong bước thứ ba, các $\lambda_{j,k+1}$ là hệ số đã được tính, không
phải các hàm cần lấy tiếp đạo hàm theo $\boldsymbol{\theta}$. Nếu không
giữ cố định, ta sẽ tối ưu một biểu thức có građien chứa thêm
\begin{equation}
 \nabla_{\boldsymbol{\theta}}(\lambda_j(\boldsymbol{\theta})J_j)
 =\lambda_j\nabla J_j+J_j\nabla\lambda_j,
 \label{p4:eq:weight-derivative}
\end{equation}
khác với thuật toán trên. Khi dùng L-BFGS với hàm mục tiêu xác định,
cần giữ tập điểm và trọng số cố định trong mỗi lượt giải thích hợp;
nếu thay chúng, các cặp sai khác građien cũ có thể không còn mô tả
cùng một hàm. Đây là yêu cầu nhất quán của cài đặt đang dùng, không
phải khẳng định không tồn tại biến thể L-BFGS ngẫu nhiên.

Cuối cùng, với $\boldsymbol{\theta}$ cố định,
$\partial J/\partial\lambda_j=J_j\geq0$.
Do đó cực tiểu hóa đồng thời một tổng phạt tuyến tính theo các trọng
số không âm, mà không thêm ràng buộc hay cơ chế khác, cho phép giảm
mục tiêu bằng cách đưa trọng số về không. Nhận xét này không loại
trừ các phương pháp học trọng số có ràng buộc, có số hạng bổ sung hoặc
có cấu trúc cực tiểu--cực đại.

\subsection{Kiến trúc có cổng và phạm vi tác động}

Một bổ sung trong công trình của Wang và cộng sự là dùng hai biểu
diễn của đầu vào ở nhiều lớp~\cite{wang2020}. Với đầu vào cột
$\mathbf{z}\in\mathbb{R}^{d+1}$ và bề rộng $h$, có thể viết
\begin{align}
 \mathbf{U}&=\sigma(\mathbf{W}_U\mathbf{z}+\mathbf{b}_U),&
 \mathbf{V}&=\sigma(\mathbf{W}_V\mathbf{z}+\mathbf{b}_V),\notag\\
 \mathbf{H}^{[1]}&=\sigma(\mathbf{W}_H\mathbf{z}+\mathbf{b}_H),\notag\\
 \mathbf{Z}^{[l]}&=\sigma(\mathbf{W}_Z^{[l]}\mathbf{H}^{[l]}
                                      +\mathbf{b}_Z^{[l]}),\notag\\
 \mathbf{H}^{[l+1]}&=(\mathbf{1}-\mathbf{Z}^{[l]})\odot\mathbf{U}
                             +\mathbf{Z}^{[l]}\odot\mathbf{V},
                      &&l=1,\ldots,L_g,\notag\\
 u_{\boldsymbol{\theta}}&=\mathbf{w}_{o}^{T}\mathbf{H}^{[L_g+1]}+b_o.
 \label{p4:eq:gated}
\end{align}
Trong đó, $\mathbf{U},\mathbf{V},\mathbf{H}^{[l]},\mathbf{Z}^{[l]}
\in\mathbb{R}^{h}$; $\mathbf{W}_U,\mathbf{W}_V,\mathbf{W}_H$
có kích thước $h\times(d+1)$;
$\mathbf{W}_Z^{[l]}\in\mathbb{R}^{h\times h}$;
$\mathbf{w}_o\in\mathbb{R}^{h}$; $L_g$ là số khối trộn;
$\odot$ là tích theo từng thành phần; $\mathbf{1}$ là vectơ toàn một.

Với $\sigma=\tanh$, các tọa độ của cổng nằm trong $(-1,1)$, nên
biểu thức trộn không nhất thiết là tổ hợp lồi. Giữ đầu vào cố định,
đạo hàm một khối theo trạng thái trước là
\begin{equation}
 \frac{\partial\mathbf{H}^{[l+1]}}{\partial\mathbf{H}^{[l]}}
 =\operatorname{diag}(\mathbf{V}-\mathbf{U})
  \operatorname{diag}\!\left(\sigma'(\mathbf{W}_Z^{[l]}
                     \mathbf{H}^{[l]}+\mathbf{b}_Z^{[l]})\right)
  \mathbf{W}_Z^{[l]}.
 \label{p4:eq:gated-jacobian}
\end{equation}
Ký hiệu $\operatorname{diag}(\mathbf{v})$ là ma trận đường chéo chứa
tọa độ của $\mathbf{v}$. Công thức xuất phát từ
$\mathbf{H}^{[l+1]}=\mathbf{U}+\mathbf{Z}^{[l]}\odot(\mathbf{V}-\mathbf{U})$
và quy tắc dây chuyền. Các đường nối với $\mathbf{U},\mathbf{V}$
tạo thêm đường truyền đạo hàm, nhưng tích các Jacobi vẫn có thể suy
giảm hoặc tăng mạnh. Kiến trúc này thay đổi cả tham số hóa và lớp
hàm biểu diễn; không có căn cứ để nói nó chỉ thay đổi tối ưu mà giữ
nguyên hoàn toàn khả năng xấp xỉ.

\section{Thiên kiến phổ và các nghiệm có thang độ dài nhỏ}
\label{sec:spectral-bias}

\subsection{Khái niệm và quy ước Fourier}

Thiên kiến phổ là xu hướng một mô hình, trong một chế độ huấn luyện
nhất định, khớp các thành phần tần số thấp sớm hơn các thành phần tần
số cao. Rahaman và cộng sự phân tích cấu trúc Fourier của mạng ReLU
và đưa ra các thí nghiệm về thứ tự học tần số~\cite{rahaman2019}.
Đây không phải khẳng định mạng không thể biểu diễn tần số cao, cũng
không phải một định lý rằng mọi kiến trúc và mọi thuật toán đều học
theo cùng thứ tự.

Với $v\in L^1(\mathbb{R}^{d})$, quy ước biến đổi Fourier là
\begin{equation}
 \widehat v(\boldsymbol{\xi})
 =\int_{\mathbb{R}^{d}}v(\mathbf{x})
                  \exp(-2\pi\mathrm{i}\boldsymbol{\xi}\cdot\mathbf{x})
                  \,\mathrm{d}\mathbf{x}.
 \label{p4:eq:fourier-transform}
\end{equation}
Trong đó, $\mathrm{i}^2=-1$ và $\boldsymbol{\xi}$ là tần số tính theo
số chu kỳ trên một đơn vị độ dài. Trên đoạn tuần hoàn có chiều dài
$\ell$, dùng hệ số
\begin{equation}
 v_k=\frac{1}{\ell}\int_0^{\ell}v(x)
                \exp(-2\pi\mathrm{i}kx/\ell)\,\mathrm{d}x,
 \qquad k\in\mathbb{Z}.
 \label{p4:eq:fourier-series}
\end{equation}
Ở đây $k$ là chỉ số nguyên của họ hàm cơ sở; tần số không gian là
$k/\ell$ và tần số góc là $\omega_k=2\pi k/\ell$.
Không đồng nhất ba đại lượng này khi ước lượng bề dày lớp biên.

Nếu các đạo hàm và điều kiện suy giảm ở vô cùng thích hợp, tích phân
từng phần cho
$\widehat{\partial_{x_j}v}=2\pi\mathrm{i}\xi_j\widehat v$.
Trên miền hữu hạn, công thức tương ứng cần điều kiện tuần hoàn hoặc
xử lý số hạng biên. Chẳng hạn, hàm hằng $1$ trên $(0,1)$ được mở rộng
bằng không có biến đổi
\begin{equation}
 \widehat{\mathbf{1}_{(0,1)}}(\xi)
 =\frac{1-\exp(-2\pi\mathrm{i}\xi)}{2\pi\mathrm{i}\xi},
 \qquad \xi\ne0.
 \label{p4:eq:window-counterexample}
\end{equation}
Sự cắt cụt ở biên tạo suy giảm cỡ $1/|\xi|$, dù hàm hằng ở trong
miền là trơn vô hạn. Do đó, phát biểu về phổ phải chỉ rõ miền và cách
xử lý biên.

\subsection{Phổ của mạng ReLU: một kết quả một chiều có chứng minh}

Mạng ReLU hữu hạn là hàm afin liên tục từng khúc. Trên một miền
kích hoạt, mỗi ReLU được thay bằng ánh xạ không hoặc ánh xạ đồng
nhất; hợp của các ánh xạ afin vẫn afin. Chia tiếp theo các mẫu
kích hoạt cho một phân hoạch hữu hạn. Trong nhiều chiều, các miền
này là các miền đa diện, và cấu trúc phổ phụ thuộc hướng của các mặt.

Để có một định lý tường minh không che giấu các điều kiện biên, xét
trường hợp một chiều có giá compact.

\begin{proposition}[Suy giảm phổ qua bước nhảy của hệ số góc]
Cho $v:\mathbb{R}\to\mathbb{R}$ liên tục, afin từng khúc và có giá
compact. Gọi $\tau_1,\ldots,\tau_M$ là các điểm gãy và
$c_j=v'(\tau_j^+)-v'(\tau_j^-)$ là bước nhảy của hệ số góc.
Với $\xi\ne0$,
\begin{equation}
 \widehat v(\xi)
 =\frac{\sum_{j=1}^{M}c_j\exp(-2\pi\mathrm{i}\xi\tau_j)}
                         {(2\pi\mathrm{i}\xi)^2},
 \qquad
 |\widehat v(\xi)|\leq
              \frac{\sum_{j=1}^{M}|c_j|}{4\pi^2\xi^2}.
 \label{p4:eq:relu-1d}
\end{equation}
\end{proposition}
Trong đó, $v'(\tau_j^\pm)$ là các giới hạn một phía của hệ số góc.
Giá compact có nghĩa $v$ bằng không ngoài một tập bị chặn.

\begin{proof}
Đạo hàm $v'$ là hàm hằng từng khúc. Với mọi hàm thử trơn
$\varphi$ có giá compact, tích phân từng phần trên từng khoảng cho
\[
 \langle v'',\varphi\rangle
 =-\int_{\mathbb{R}}v'(x)\varphi'(x)\,\mathrm{d}x
 =\sum_{j=1}^{M}c_j\varphi(\tau_j).
\]
Các đầu mút bên trong đóng góp chênh lệch hệ số góc, còn số hạng ở
vô cùng bằng không. Vì vậy $v''=\sum_jc_j\delta_{\tau_j}$ theo nghĩa
phân bố, với $\delta_{\tau_j}$ là khối lượng Dirac tại $\tau_j$.
Lấy biến đổi Fourier, dùng
$\widehat{\delta_{\tau_j}}(\xi)=\exp(-2\pi\mathrm{i}\xi\tau_j)$:
\[
 (2\pi\mathrm{i}\xi)^2\widehat v(\xi)
 =\sum_jc_j\exp(-2\pi\mathrm{i}\xi\tau_j).
\]
Chia cho $(2\pi\mathrm{i}\xi)^2$ và dùng bất đẳng thức tam giác cho
\eqref{p4:eq:relu-1d}.
\end{proof}

Nếu $v$ có hằng số Lipschitz $L_v$, mỗi $|c_j|\leq2L_v$.
Với mạng ReLU có cấu trúc như \eqref{p4:eq:network}, tính chất
Lipschitz của hợp hàm còn cho
\begin{equation}
 L_v\leq\prod_{l=1}^{D}\|\mathbf{W}^{[l]}\|_2,
 \label{p4:eq:lip-network}
\end{equation}
trong đó chuẩn ma trận là chuẩn phổ. Thật vậy, mỗi ánh xạ afin có
hằng số Lipschitz bằng chuẩn phổ của ma trận trọng số, ReLU có hằng
số bằng một và các hằng số được nhân khi hợp hàm.
Nếu mạng thuộc trường hợp có giá compact của mệnh đề và
$|\widehat v(\xi)|\geq A>0$, thì
\begin{equation}
 \prod_{l=1}^{D}\|\mathbf{W}^{[l]}\|_2
 \geq L_v\geq\frac{2\pi^2A\xi^2}{M}.
 \label{p4:eq:spectral-parameter-cost}
\end{equation}
Đây là điều kiện cần có giả thiết, với $M$ là số điểm gãy thực tế;
việc thay đổi kiến trúc cũng có thể thay đổi $M$. Nó không cho một
cận phổ dụng chỉ theo chuẩn Euclid của toàn bộ tham số.

\paragraph{Điểm cần giữ đúng khi đọc kết quả nhiều chiều.}
Phân tích của Rahaman và cộng sự chỉ ra sự suy giảm không đẳng hướng:
trong bối cảnh của bài báo, các hướng tổng quát có tốc độ cỡ
$\|\boldsymbol{\xi}\|^{-(d+1)}$, nhưng ở các hướng đặc biệt liên quan
đến pháp tuyến của mặt phân hoạch, tốc độ có thể chậm tới cỡ
$\|\boldsymbol{\xi}\|^{-2}$~\cite[Mục~2]{rahaman2019}.
Vì vậy không viết một cận với số mũ $d+1$ cho mọi hướng chỉ từ giả
thiết mạng ReLU được xét trên miền bị chặn. Cũng không suy ra kết
quả đó bằng cách khẳng định mọi đạo hàm bậc $d+1$ đều là tổng hữu hạn
các khối lượng Dirac: đạo hàm phân bố trong nhiều chiều còn chứa các
độ đo trên mặt và các đạo hàm của phân bố.

\subsection{Tính trơn của hàm kích hoạt và giới hạn của lập luận phổ}

Với $v$ tuần hoàn đủ trơn, tích phân từng phần $m$ lần trong
\eqref{p4:eq:fourier-series}, khi các đạo hàm đến bậc $m-1$ khớp ở
hai đầu mút, cho
\begin{equation}
 v_k=\frac{1}{\ell(\mathrm{i}\omega_k)^m}
       \int_0^{\ell}v^{(m)}(x)\exp(-\mathrm{i}\omega_kx)\,\mathrm{d}x,
 \qquad
 |v_k|\leq\frac{\|v^{(m)}\|_{L^1(0,\ell)}}
                              {\ell|\omega_k|^m},\quad k\ne0.
 \label{p4:eq:smooth-decay}
\end{equation}
Ở đây $m$ là số lần tích phân từng phần. Tốc độ suy giảm đi kèm một
hằng số phụ thuộc đạo hàm bậc $m$. Khi trọng số mạng tăng hoặc khi
mạng biểu diễn một cấu trúc dốc, hằng số này cũng có thể lớn.
Tính trơn của $\tanh$ không tự nó cho một cận phổ đồng đều theo tất
cả bộ trọng số, đặc biệt nếu dữ kiện không tuần hoàn và bị cắt cụt
tại biên. Do đó không chuyển nguyên số mũ và hằng số của định lý cho
ReLU sang mạng dùng $\tanh$.

Một hàm có phổ suy giảm ở tần số lớn vẫn có thể chứa một thành phần
tần số cao đáng kể. Ví dụ, $v(x)=\sin(2\pi Kx)$ là hàm giải tích
nhưng mang tần số $K$, với $K$ có thể lớn tùy ý. Tính trơn là thuộc
tính của từng hàm, còn thiên kiến phổ liên quan thêm đến cách tham
số hóa, khởi tạo và quỹ đạo huấn luyện.

\subsection{Từ phổ của hàm đến động lực học huấn luyện}

Để phân biệt hai vấn đề, xét trước bài toán hồi quy trên các điểm cố
định $\mathbf{z}_1,\ldots,\mathbf{z}_M$ với
\begin{equation}
 J_{\mathrm{reg}}(\boldsymbol{\theta})
 =\frac{1}{2M}\|\mathbf{F}(\boldsymbol{\theta})-\mathbf{y}\|_2^2,
 \quad F_i(\boldsymbol{\theta})=u_{\boldsymbol{\theta}}(\mathbf{z}_i).
 \label{p4:eq:regression}
\end{equation}
Trong đó, $\mathbf{y}\in\mathbb{R}^{M}$ chứa giá trị đích;
$\mathbf{e}=\mathbf{F}-\mathbf{y}$ là vectơ sai số;
$\mathbf{B}=D_{\boldsymbol{\theta}}\mathbf{F}$ có kích thước
$M\times P$. Dòng građien cho
\begin{equation}
 \dot{\boldsymbol{\theta}}=-\frac{1}{M}\mathbf{B}^{T}\mathbf{e},
 \qquad
 \dot{\mathbf{e}}=\mathbf{B}\dot{\boldsymbol{\theta}}
 =-\mathbf{K}(\boldsymbol{\theta})\mathbf{e},
 \qquad \mathbf{K}:=\frac{1}{M}\mathbf{B}\mathbf{B}^{T}.
 \label{p4:eq:kernel-flow}
\end{equation}
Ma trận $\mathbf{K}$ là ma trận của nhân tiếp tuyến nơ-ron trên tập
điểm; nguồn gốc của cách phân tích này được trình bày trong
\cite{jacot2018}. Đẳng thức \eqref{p4:eq:kernel-flow} là hệ quả trực
tiếp của quy tắc dây chuyền, chưa cần giả thiết bề rộng vô hạn.

Nếu trong một xấp xỉ tuyến tính hóa ta giữ
$\mathbf{K}(\boldsymbol{\theta})\approx\mathbf{K}_0$ không đổi và
$\mathbf{K}_0\mathbf{v}_j=\kappa_j\mathbf{v}_j$, thì
\begin{equation}
 \mathbf{v}_j^T\mathbf{e}(s)
 =\exp(-\kappa_js)\,\mathbf{v}_j^T\mathbf{e}(0).
 \label{p4:eq:kernel-modes}
\end{equation}
Các $\kappa_j\geq0$ là trị riêng. Hướng có trị riêng nhỏ giảm chậm;
hướng có trị riêng bằng không không giảm trong mô hình đóng băng.
Tuy nhiên, vectơ riêng của một ma trận nhân trên các điểm tùy ý
không tự động là các hàm Fourier, và việc sắp trị riêng giảm dần
không đồng nghĩa sắp tần số tăng dần.

Để có kết luận theo tần số, cần thêm cấu trúc. Chẳng hạn, trên miền
tuần hoàn, nếu toán tử nhân là tích chập
$\mathcal{K}h(x)=\ell^{-1}\int_0^\ell K(x-y)h(y)\,\mathrm{d}y$,
thì các hàm $\exp(\mathrm{i}\omega_kx)$ là hàm riêng. Gọi $\kappa_k$
là trị riêng tương ứng. Với dòng $\partial_se=-\mathcal{K}e$,
\begin{equation}
 e_k(s)=\exp(-\kappa_ks)e_k(0).
 \label{p4:eq:fourier-training}
\end{equation}
Chỉ khi các $\kappa_k$ nhỏ ở tần số cao mới suy ra những thành phần
đó được học chậm từ công thức này. Đây là một cơ chế có điều kiện,
không phải một tính chất đã chứng minh cho mọi mạng hữu hạn.

\subsection{Tác động của toán tử vi phân lên phổ sai số}

Trên miền tuần hoàn, xét toán tử tuyến tính hệ số hằng
\begin{equation}
 \mathcal{L}=\sum_{|\boldsymbol{\alpha}|\leq m}
       c_{\boldsymbol{\alpha}}\partial^{\boldsymbol{\alpha}},
 \qquad
 p(\mathbf{z})=\sum_{|\boldsymbol{\alpha}|\leq m}
       c_{\boldsymbol{\alpha}}\mathbf{z}^{\boldsymbol{\alpha}}.
 \label{p4:eq:symbol}
\end{equation}
Trong đó, $\boldsymbol{\alpha}\in\mathbb{N}_0^d$ là đa chỉ số,
$|\boldsymbol{\alpha}|$ là tổng các thành phần,
$\partial^{\boldsymbol{\alpha}}$ là đạo hàm hỗn hợp và $p$ là đa thức
biểu trưng. Nếu $\mathcal{L}u=f$, $e=v-u$, $r=\mathcal{L}v-f$,
thì tính tuyến tính cho $r=\mathcal{L}e$. Với tần số góc
$\boldsymbol{\omega}_k$, quy tắc đạo hàm Fourier cho
\begin{equation}
 r_k=p(\mathrm{i}\boldsymbol{\omega}_k)e_k.
 \label{p4:eq:residual-symbol}
\end{equation}
Đây là đẳng thức chính xác trong bối cảnh đã nêu. Để thay
$|p(\mathrm{i}\boldsymbol{\omega})|$ bằng một đại lượng tương đương
$\|\boldsymbol{\omega}\|^m$ đồng đều theo mọi hướng, cần thêm điều
kiện không suy biến của biểu trưng chính, chẳng hạn tính elliptic.
Phản ví dụ là $\mathcal{L}=\partial_{x_1}$: biểu trưng bằng không
trên các tần số $(0,\omega_2)$ dù chuẩn tần số lớn tùy ý.

Với Poisson một chiều, $\mathcal{L}=-\partial_{xx}$ nên
$r_k=\omega_k^2e_k$. Parseval cho
\begin{equation}
 \frac{1}{\ell}\|r\|_{L^2(0,\ell)}^2
 =\sum_{k\in\mathbb{Z}}\omega_k^4|e_k|^2.
 \label{p4:eq:poisson-spectral-loss}
\end{equation}
Ở đây, chuẩn phần dư phạt sai số ở tần số cao mạnh hơn theo lũy thừa
bốn. Thành phần có tần số bằng không bị toán tử bỏ qua; trong bài toán Poisson tuần hoàn
cần thêm điều kiện xác định giá trị trung bình của nghiệm.

Để thấy hệ quả về điều kiện số mà không lẫn với tham số hóa mạng,
chọn một họ hữu hạn các hàm sin trực chuẩn $\varphi_1,\ldots,\varphi_K$
với $-\varphi_j''=\omega_j^2\varphi_j$, $\omega_j=2\pi j/\ell$.
Viết $e=\sum_{j=1}^{K}a_j\varphi_j$. Khi đó
\begin{equation}
 \tfrac12\|\mathcal{L}e\|_{L^2}^2
 =\tfrac12\sum_{j=1}^{K}\omega_j^4a_j^2,
 \qquad
 \nabla_{\mathbf{a}}^2J=\operatorname{diag}(\omega_1^4,\ldots,\omega_K^4).
 \label{p4:eq:spectral-hessian}
\end{equation}
Số điều kiện theo các hệ số $\mathbf{a}$ bằng $K^4$.
Nếu các hệ số phụ thuộc phi tuyến vào tham số mạng, Hessian tham số
còn chứa Jacobi và đạo hàm bậc hai của ánh xạ đó; không được đồng
nhất nó với ma trận đường chéo ở \eqref{p4:eq:spectral-hessian}.

Tác động của đạo hàm cũng không chỉ theo một chiều bất lợi. Trong
mô hình nhân cố định, tuần hoàn và bất biến tịnh tiến, với mất mát
$\tfrac12\|\mathcal{L}e\|^2$ theo chuẩn đã chuẩn hóa, dòng sai số là
$\partial_se=-\mathcal{K}\mathcal{L}^{*}\mathcal{L}e$.
Do đó
\begin{equation}
 \partial_se_k=-\kappa_k|p(\mathrm{i}\boldsymbol{\omega}_k)|^2e_k.
 \label{p4:eq:pinn-kernel-rate}
\end{equation}
Ký hiệu $\mathcal{L}^{*}$ là toán tử liên hợp theo tích vô hướng
$L^2$ tuần hoàn. Công thức được suy ra vì $\mathcal{L}$ nhân mỗi hệ
số với $p$, toán tử liên hợp nhân với liên hợp phức của $p$, còn
$\mathcal{K}$ nhân với $\kappa_k$. Hệ số đạo hàm có thể tăng tốc độ
giảm liên tục của một số tần số, đồng thời làm tăng trị riêng lớn
nhất và thu hẹp bước ổn định của thuật toán rời rạc. Kết quả sau cùng
phụ thuộc tích $\kappa_k|p|^2$, không chỉ phụ thuộc riêng độ trơn của
mạng hoặc bậc đạo hàm.

\subsection{Lớp biên, lớp chuyển tiếp và phương trình Burgers}

Một ví dụ chính xác về lớp biên là
\begin{equation}
 -\nu u''+u'=1,\quad 0<x<1,\quad u(0)=u(1)=0,
 \qquad
 u_\nu(x)=x-\frac{\exp(x/\nu)-1}{\exp(1/\nu)-1}.
 \label{p4:eq:boundary-layer}
\end{equation}
Để suy ra, đặt $v=u'$: phương trình $-\nu v'+v=1$ có nghiệm
$v=1+C\exp(x/\nu)$. Tích phân và áp hai điều kiện biên cho biểu thức
trên. Khi $\nu$ nhỏ, số hạng hiệu chỉnh gần $x=1$ có dạng
$\exp(-(1-x)/\nu)$, nên biến thiên trên khoảng cách cỡ $\nu$.
Không phải mọi vế phải và mọi dữ kiện biên đều tạo lớp biên không
tầm thường; ở đây chúng đã được chỉ rõ.

Với Burgers có độ nhớt trên toàn trục, một hồ sơ chuyển tiếp dừng là
\begin{equation}
 U(x)=-A\tanh\!\left(\frac{A(x-x_0)}{2\nu}\right),
 \qquad A>0.
 \label{p4:eq:burgers-profile}
\end{equation}
Trong đó, $A$ là biên độ hai trạng thái giới hạn và $x_0$ là tâm lớp.
Thật vậy, phương trình $UU'-\nu U''=0$ tích phân một lần cho
$\nu U'=(U^2-A^2)/2$ nếu $U\to\pm A$ và $U'\to0$ ở hai phía.
Thay \eqref{p4:eq:burgers-profile} vào cho
\[
 U'=-\frac{A^2}{2\nu}
       \operatorname{sech}^{2}\!\left(\frac{A(x-x_0)}{2\nu}\right),
 \qquad \max_x|U'|=\frac{A^2}{2\nu},
\]
với $\operatorname{sech}z=1/\cosh z$. Nếu định nghĩa bề dày bằng tỉ
số độ chênh hai trạng thái và độ dốc lớn nhất thì
\begin{equation}
 \delta:=\frac{2A}{\max|U'|}=\frac{4\nu}{A}.
 \label{p4:eq:layer-width}
\end{equation}
Đây là kết quả của hồ sơ chính xác vừa xét, không phải số đo của
bài toán giá trị đầu $u(x,0)=-\sin(\pi x)$ trên $[-1,1]$.
Với $A=1$ và $\nu=0{,}01/\pi$, công thức cho
$\delta\approx1{,}2732\times10^{-2}$ như một thang tham chiếu.
Không dùng nó để gán thời điểm hoặc độ dốc cực đại cho một thí nghiệm
chưa được thực hiện.

Mối liên hệ giữa thang không gian và thang phổ có thể suy ra chính
xác từ phép đổi biến. Với $F\in L^1(\mathbb{R})$ và
$F_\delta(x)=F((x-x_0)/\delta)$,
\begin{equation}
 \widehat{F_\delta}(\xi)
 =\delta\exp(-2\pi\mathrm{i}\xi x_0)\widehat F(\delta\xi).
 \label{p4:eq:scaling-fourier}
\end{equation}
Đặt $y=(x-x_0)/\delta$ trong tích phân Fourier là đủ để thu được
công thức. Khi $\delta$ giảm, phổ theo biến $\xi$ giãn ra theo
$1/\delta$. Có thể áp dụng lập luận này cho đạo hàm cục bộ hóa của
hồ sơ chuyển tiếp, vì bản thân $U$ trong \eqref{p4:eq:burgers-profile}
không thuộc $L^1(\mathbb{R})$.
Như vậy thang tần số cao xuất hiện một cách có cơ sở, nhưng không
xác định một chỉ số Fourier duy nhất nếu chưa nêu chu kỳ và ngưỡng
cắt phổ.

Với nghiệm Burgers tổng quát, quan hệ phần dư--sai số còn là phi
tuyến. Đặt $v=u+e$, với $u_t+uu_x-\nu u_{xx}=0$. Khai triển
$(u+e)(u_x+e_x)$ cho
\begin{equation}
 r(v)=e_t+u e_x+u_xe+e e_x-\nu e_{xx}.
 \label{p4:eq:burgers-error}
\end{equation}
Ngay cả khi bỏ số hạng bậc hai $e e_x$ để tuyến tính hóa, các hệ số
$u,u_x$ vẫn phụ thuộc $(x,t)$. Phép biến đổi Fourier của phép nhân
với hệ số biến thiên tạo tích chập giữa các tần số; không còn là
phép nhân với một đa thức duy nhất như
\eqref{p4:eq:residual-symbol}. Phân tích biểu trưng hệ số hằng chỉ là
một mô hình giải thích cục bộ nếu các giả thiết đó không thỏa.

\subsection{Hệ quả đối với thiết kế và đánh giá PINN}

Các phân tích trên gợi ý phân biệt bốn câu hỏi: lớp hàm có xấp xỉ
được nghiệm trong chuẩn cần thiết không; tập phối trí có đánh giá
đúng phần dư không; thuật toán có giảm đủ hàm mục tiêu không; và
chuẩn phần dư đó có kiểm soát sai số nghiệm không. Những câu hỏi này
liên hệ với nhau, nhưng không thể thay thế cho nhau.

Một cách thay đổi biểu diễn đầu vào là dùng các đặc trưng Fourier
\begin{equation}
 \gamma(\mathbf{z})=
 \begin{bmatrix}
 \cos(2\pi\mathbf{B}\mathbf{z})\\
 \sin(2\pi\mathbf{B}\mathbf{z})
 \end{bmatrix},\qquad
 \mathbf{B}\in\mathbb{R}^{m_f\times(d+1)}.
 \label{p4:eq:fourier-features}
\end{equation}
Trong đó, các hàng của $\mathbf{B}$ là các vectơ tần số được chọn,
các hàm lượng giác tác động theo tọa độ và $m_f$ là số tần số.
Mạng nhận $\gamma(\mathbf{z})$ thay cho hoặc cùng với đầu vào ban
đầu. Cơ chế này được nghiên cứu trong \cite{tancik2020}.
Nó đưa các thang dao động vào biểu diễn, nhưng khi tính phần dư phải
lấy cả đạo hàm của $\gamma$; các hệ số tần số lớn cũng xuất hiện
trong các đạo hàm đó.

\begin{table}[htbp]
\centering\small
\caption{Cơ chế tác động và giới hạn cần kiểm tra của một số lựa chọn.}
\label{p4:tab:remedies}
\begin{tabularx}{\textwidth}{@{}>{\raggedright\arraybackslash}p{0.23\textwidth}>{\raggedright\arraybackslash}X>{\raggedright\arraybackslash}X@{}}
\toprule
\textbf{Lựa chọn}&\textbf{Tác động chính}&\textbf{Điều cần kiểm tra}\\
\midrule
Khử thứ nguyên và điều chỉnh trọng số&
Thay đổi tỉ lệ các thành phần và hướng cập nhật.&
Giá trị từng mất mát, độ lớn và góc giữa các građien.\\
\addlinespace
Đặc trưng Fourier hoặc đổi thang đầu vào&
Thay đổi biểu diễn và độ nhạy theo tham số.&
Dải tần được đưa vào, đạo hàm bậc cao và sai số trên điểm mới.\\
\addlinespace
Lấy mẫu thích nghi&
Tăng mật độ đánh giá ở vùng được chọn.&
Chuẩn mục tiêu có thay đổi không; trọng số lấy mẫu có cần hiệu chỉnh không.\\
\addlinespace
Ràng buộc cứng&
Áp chính xác một số điều kiện theo cấu trúc.&
Độ trơn của hàm nâng và nhân tử; tính tương thích và khả năng xấp xỉ.\\
\addlinespace
Chia miền hoặc chia khoảng thời gian&
Thay đổi kích thước và thang của các bài toán con.&
Điều kiện ghép nối, sự truyền sai số và chi phí tổng.\\
\bottomrule
\end{tabularx}
\end{table}

Không có kết luận chung rằng chính quy hóa luôn bất lợi cho lớp
biên: giảm chuẩn trọng số có thể làm khó việc khớp một số dao động,
nhưng cũng có thể giảm nhạy cảm với nhiễu và quá khớp tập phối trí.
Tương tự, một kỹ thuật thay đổi kiến trúc hoặc thang đo có thể đồng
thời tác động lên xấp xỉ và tối ưu; không nên xem các tầng phân tích
là các cơ chế hoàn toàn độc lập.

Trong đánh giá số, cần báo cáo riêng các thành phần mất mát, kích
thước và quy tắc lấy mẫu, chiều dài bước, trọng số, kiến trúc và
chuẩn sai số. Khi có nghiệm tham chiếu, sai số tương đối theo $L^2$
là $\|u_{\boldsymbol{\theta}}-u_{\mathrm{ref}}\|_{L^2}/
\|u_{\mathrm{ref}}\|_{L^2}$, với mẫu số khác không; sai số theo chuẩn
cực đại là một đại lượng khác. Đối với lớp biên hoặc lớp chuyển tiếp,
nên kiểm tra thêm vùng biến thiên nhanh trên tập điểm đủ mịn. Các
chỉ tiêu này giúp phân biệt nghiệm có phần dư nhỏ tại điểm huấn
luyện với nghiệm thực sự chính xác trong chuẩn cần quan tâm.
